\begin{abstract}
\noindent This study presents a methodology for detecting, analysing, and classifying Marine heatwaves (MHWs) at hourly resolution. 
MHWs are traditionally identified from daily Sea Surface Temperature (SST)
records, yet this framework does not transfer directly to higher temporal resolution. At hourly scale,
SST reveals sub-daily variability that is largely smoothed out in daily means, including intermittent
threshold exceedances, short-lived interruptions, diurnal cycling, and nocturnal heat retention. These
features complicate the definition of event continuity and persistence, but also provide an opportunity
to examine the internal structure of marine heat extremes in substantially greater detail.

\noindent 
Events are identified through a duration--gap framework that preserves short interruptions while
maintaining event continuity, and are subsequently decomposed into individual diurnal cycles. Their
intra-event morphology is used to construct an hourly-scale taxonomy that characterises MHWs not only
in terms of conventional duration and intensity, but also through persistence, interruption structure,
cycle-to-cycle evolution, diurnal heating, nocturnal retention, and other sub-daily dynamics that are
not accessible from daily observations. To enable the methodology to scale from individual time
series to multi-year gridded SST datasets, a GPU-accelerated implementation, \texttt{CUDAHeat}, is
developed using CUDA-enabled array computation. The computational architecture preserves the complete
temporal record of each spatial point while partitioning only the spatial domain into dynamically sized
batches according to available GPU memory. VRAM-aware batch adaptation, automatic out-of-memory
recovery, and explicit support for both 32-bit and 64-bit floating-point computation allow the same
workflow to operate across different GPU memory capacities without altering the temporal definition
of an event.

\noindent Benchmarking on a hourly gridded SST dataset comprising approximately 1.16 billion values
shows a $23.4\times$ acceleration of the core computation using GPU 32-bit arithmetic relative to the
CPU 64-bit reference, with a $15.9\times$ speedup for the complete processing pipeline. Scientific
validation yields approximately $99.9991\%$ agreement in threshold-exceedance classification across
tens of millions of comparisons, with the remaining differences confined predominantly to values
extremely close to the detection threshold. These results demonstrate that the hourly framework can
retain its temporal and taxonomic detail while remaining computationally tractable for substantially
larger spatial datasets. This combination of an hourly event taxonomy with scalable GPU computation provides a framework for
investigating marine heat extremes from their conventional daily representation down to their
sub-daily structure, and establishes a basis for subsequent analysis of coupled marine--atmospheric
heat extremes and coastal thermal stress.
\end{abstract}